\section{Introdução e objetivos}
\label{sec:objetivos}

Um amplificador operacional real se comporta de forma um pouco diferente
do modelo ideal. Entram pequenas correntes contínuas
pelos terminais de entrada (as correntes de polarização $I_{B+}$ e
$I_{B-}$, cuja diferença é a corrente de offset $I_{OS}$). Existe uma
tensão de offset de entrada $V_{OS}$. O ganho diferencial $A_{vd}$ é
grande, mas finito. Sinais aplicados igualmente nas duas entradas (modo
comum) não são totalmente rejeitados, o que é medido pela CMRR. Além
disso, a saída não consegue variar mais rápido que um certo limite,
chamado de slew rate. No OP07, que é um amplificador de precisão, esses efeitos são muito pequenos, então cada um
precisa de um circuito de teste próprio para virar uma medida possível.

O roteiro usa o Micro-Cap. Aqui as simulações foram feitas no LTspice, rodando em containers Docker do repositório
\texttt{claw-spice}, com um container independente para cada simulação.
Os objetivos, seguindo os itens do roteiro, são:
medir $I_{B+}$ e $I_{B-}$ e calcular $I_{OS}$; levantar a curva de
transferência em malha aberta e identificar $V_{OS}$; medir $A_{vd}$ e
$A_{cm}$ pela derivada das curvas e calcular a CMRR; extrair o slew rate
da resposta ao degrau; e comparar todos os valores com a folha de dados
do fabricante e com o amplificador ideal.
